\subsection{Quaternion Specialization of the Two-by-Two Map}
The intrinsic quaternion basis from Eq.\ \eqref{eq:quaternion-basis} is
obtained by restricting the same Pauli matrix map already defined in
Eq.\ \eqref{eq:matrix-paravector-map}. Its three basis values are
\begin{equation}
\begin{aligned}
\matf{\pv{q}_1}
&=-i\mat{\sigma}_1
=\begin{bmatrix}0&-i\\-i&0\end{bmatrix},&
\matf{\pv{q}_2}
&=-i\mat{\sigma}_2
=\begin{bmatrix}0&-1\\1&0\end{bmatrix},\\[0.35em]
\matf{\pv{q}_3}
&=-i\mat{\sigma}_3
=\begin{bmatrix}-i&0\\0&i\end{bmatrix}.
\end{aligned}
\label{eq:quaternion-matrices}
\end{equation}
For the quaternion in Eq.\ \eqref{eq:quaternion-paravector}, the two-by-two
value of the map is
\begin{equation}
\matf{\pv{Q}}
=
\begin{bmatrix}
a-iB_3&-B_2-iB_1\\
B_2-iB_1&a+iB_3
\end{bmatrix}.
\label{eq:quaternion-matrix-map}
\end{equation}
Because intrinsic star fixes \(\pv{Q}\), its induced matrix operation also
fixes \(\matf{\pv{Q}}\):
\begin{equation*}
\conj{\matf{\pv{Q}}}
=\matf{\conj{\pv{Q}}}
=\matf{\pv{Q}}.
\end{equation*}
This star is the induced operation from
Eq.\ \eqref{eq:matrix-operation-dictionary}; it is not entrywise matrix
\(\mathsf C\), which does not fix the displayed matrix in general.
The Hermitian and adjugate operations coincide on this specialization:
\begin{equation*}
\herm{\matf{\pv{Q}}}
=\matf{\herm{\pv{Q}}}
=\adjf{\matf{\pv{Q}}}.
\end{equation*}
Consequently,
\begin{equation}
\begin{aligned}
\matf{\pv{Q}}\herm{\matf{\pv{Q}}}
&=\herm{\matf{\pv{Q}}}\matf{\pv{Q}}
=(a^2+\vct{B}^2)\I,\\
\detf{\matf{\pv{Q}}}&=a^2+\vct{B}^2,\qquad
\eqfunc{tr}(\matf{\pv{Q}})=2a.
\end{aligned}
\label{eq:quaternion-matrix-norm}
\end{equation}
Thus a unit quaternion is represented by a matrix whose Hermitian product is
\(\I\) and whose determinant is \(1\). This is a specialization of the
single Pauli matrix realization, not a second representation of the
paravector algebra.
